% GENERATED FILE. Edit canonical lessons, then run npm run generate:books.
% canonical-source-hash: fnv1a64:33b292c1f8a92dfd
% canonical-lessons: GE-C184-parken, GE-C184-pflegen, GE-C184-pfluecken, GE-C184-pruefen, GE-C184-reinigen

\chapter{To park, To look after, To pick, To test, To clean}
\label{ch:ge-to-park-to-look-after-to-pick-to-test-to-clean}
\hlchaptermodality{184}

\begin{chapteropening}
\textbf{By the end of this chapter:} \emph{I can say to park, to look after, to pick, to test and to clean.}

Complete the last lesson of chapter 184: I can say to park, to look after, to pick, to test and to clean.
\end{chapteropening}

\section[parken]{parken — to park}

\label{lesson:GE-C184-parken}



\begin{quote}
Before the new one: say the German for to sew, then the German for to sneeze.
\end{quote}

\subsection*{You'll want to know: parken}
\textbf{parken} — \textquotedblleft{}to park\textquotedblright{}.

\begin{cousinweb}
\textbf{parken} is taken from English \emph{to park}, built on \emph{park}, which came from Old French \emph{parc}, \textquotedblleft{}enclosure\textquotedblright{} (EtymDB 2.1). A car park is an enclosure.
\end{cousinweb}

\begin{grammarlens}[title={What you've built}]
One more word for this chapter.
\end{grammarlens}

\subsection*{Your turn}
\begin{itemize}
\raggedright
  \item \emph{Say it:} \emph{parken}
  \item \emph{Say it:} \emph{parken}, once more
  \item \emph{Say it:} say \emph{parken}
  \item \emph{Recall:} say the German for to sew, then the German for to sneeze, then say \emph{parken} again
\end{itemize}

\begin{tcolorbox}[breakable,colback=teal!4,colframe=teal!35!black,title={Before you move on}]
What does parken mean? (\textquotedblleft{}To park\textquotedblright{}.) Say it once more.
\end{tcolorbox}
\section[pflegen]{pflegen — to look after, to nurse}

\label{lesson:GE-C184-pflegen}



\begin{quote}
Before the new one: say the German for to sneeze, then the German for to park.
\end{quote}

\subsection*{You'll want to know: pflegen}
\textbf{pflegen} — \textquotedblleft{}to look after, to nurse\textquotedblright{}.

\begin{cousinweb}
\textbf{pflegen} $\leftarrow$ Middle High German \emph{phlëgen} $\leftarrow$ Old High German \emph{plegan}, \textquotedblleft{}to look after, to be responsible for\textquotedblright{} (EtymDB 2.1) — inherited. English \emph{plight}, a pledge, is a cousin.
\end{cousinweb}

\begin{grammarlens}[title={What you've built}]
One more word for this chapter.
\end{grammarlens}

\subsection*{Your turn}
\begin{itemize}
\raggedright
  \item \emph{Say it:} \emph{pflegen}
  \item \emph{Say it:} \emph{pflegen}, once more
  \item \emph{Say it:} say \emph{pflegen}
  \item \emph{Recall:} say the German for to sneeze, then the German for to park, then say \emph{pflegen} again
\end{itemize}

\begin{tcolorbox}[breakable,colback=teal!4,colframe=teal!35!black,title={Before you move on}]
What does pflegen mean? (\textquotedblleft{}To look after, to nurse\textquotedblright{}.) Say it once more.
\end{tcolorbox}
\section[pflücken]{pflücken — to pick (flowers, fruit)}

\label{lesson:GE-C184-pfluecken}



\begin{quote}
Before the new one: say the German for to park, then the German for to look after.
\end{quote}

\subsection*{You'll want to know: pflücken}
\textbf{pflücken} — \textquotedblleft{}to pick (flowers, fruit)\textquotedblright{}.

\begin{cousinweb}
\textbf{pflücken} $\leftarrow$ Middle High German \emph{pflocken} — inherited; EtymDB 2.1 lists it beside English \emph{pluck} and Dutch \emph{plukken}.
\end{cousinweb}

\begin{grammarlens}[title={What you've built}]
One more word for this chapter.
\end{grammarlens}

\subsection*{Your turn}
\begin{itemize}
\raggedright
  \item \emph{Say it:} \emph{pflücken}
  \item \emph{Say it:} \emph{pflücken}, once more
  \item \emph{Say it:} say \emph{pflücken}
  \item \emph{Recall:} say the German for to park, then the German for to look after, then say \emph{pflücken} again
\end{itemize}

\begin{tcolorbox}[breakable,colback=teal!4,colframe=teal!35!black,title={Before you move on}]
What does pflücken mean? (\textquotedblleft{}To pick (flowers, fruit)\textquotedblright{}.) Say it once more.
\end{tcolorbox}
\section[prüfen]{prüfen — to test, to check}

\label{lesson:GE-C184-pruefen}



\begin{quote}
Before the new one: say the German for to look after, then the German for to pick.
\end{quote}

\subsection*{You'll want to know: prüfen}
\textbf{prüfen} — \textquotedblleft{}to test, to check\textquotedblright{}.

\begin{cousinweb}
\textbf{prüfen} $\leftarrow$ Middle High German \emph{brüeven} (EtymDB 2.1), from Old French \emph{prover} $\leftarrow$ Latin \emph{probāre}, \textquotedblleft{}to test\textquotedblright{} — borrowed, like English \emph{prove}.
\end{cousinweb}

\begin{grammarlens}[title={What you've built}]
One more word for this chapter.
\end{grammarlens}

\subsection*{Your turn}
\begin{itemize}
\raggedright
  \item \emph{Say it:} \emph{prüfen}
  \item \emph{Say it:} \emph{prüfen}, once more
  \item \emph{Say it:} say \emph{prüfen}
  \item \emph{Recall:} say the German for to look after, then the German for to pick, then say \emph{prüfen} again
\end{itemize}

\begin{tcolorbox}[breakable,colback=teal!4,colframe=teal!35!black,title={Before you move on}]
What does prüfen mean? (\textquotedblleft{}To test, to check\textquotedblright{}.) Say it once more.
\end{tcolorbox}
\section[reinigen]{reinigen — to clean}

\label{lesson:GE-C184-reinigen}



\begin{quote}
Before the new one: say the German for to pick, then the German for to test.
\end{quote}

\subsection*{You'll want to know: reinigen}
\textbf{reinigen} — \textquotedblleft{}to clean\textquotedblright{}.

\begin{cousinweb}
\textbf{reinigen} is the word for \textquotedblleft{}pure, clean\textquotedblright{}, Old High German \emph{reini} $\leftarrow$ Proto-Germanic \emph{*hrainiz} (EtymDB 2.1), + the verb ending \emph{-igen}: to make clean.
\end{cousinweb}

\begin{grammarlens}[title={What you've built}]
One more word for this chapter.
\end{grammarlens}

\subsection*{Your turn}
\begin{itemize}
\raggedright
  \item \emph{Say it:} \emph{reinigen}
  \item \emph{Say it:} \emph{reinigen}, once more
  \item \emph{Say it:} say \emph{reinigen}
  \item \emph{Recall:} say the German for to pick, then the German for to test, then say \emph{reinigen} again
\end{itemize}

\begin{tcolorbox}[breakable,colback=teal!4,colframe=teal!35!black,title={Before you move on}]
What does reinigen mean? (\textquotedblleft{}To clean\textquotedblright{}.) Say it once more.
\end{tcolorbox}
